\documentclass{report}
\usepackage{amsmath,amssymb}
\setlength{\parindent}{0mm}
\setlength{\parskip}{1em}
\begin{document}
\begin{center}
\rule{6in}{1pt} \
{\large Michele Benzi \\
{\bf Numerical solution of differential equations on graphs}}

Department of Mathematics and Computer Science \\ Emory University \\ 400 Dowman Drive \\ Atlanta \\ GA 30322 \\ USA
\\
{\tt benzi@mathcs.emory.edu}\end{center}

Diffusion-type models on graphs and networks have been widely used to
model phenomena such as consensus and belief propagation in social
networks. These simple models usually involve only combinatorial graphs,
possibly weighted, and lead to systems of ODEs involving the graph
Laplacian. More sophisticated models may be obtained by studying
evolution processes on metric graphs, and
replacing the graph Laplacian with a more general differential operator,
leading to a PDE on the graph. Now not only initial conditions but also
interface conditions at the vertices must be supplied.

In this talk I will discuss
the numerical solution of differential equations posed on graphs or networks.
More specifically, the talk is concerned with
quantum graphs, which are metric graphs endowed with a self-adjoint
differential operator (Hamiltonian)
acting on functions defined on the graph's edges with
suitable side conditions. I will describe and analyze the use of
a linear finite element method for computing steady-state
solutions first.
The solution of the discrete equations is achieved by means of a
(non-overlapping) domain decomposition approach.
For model elliptic problems and a wide class of graphs, we show that
a combination of Schur complement reduction and diagonally preconditioned
conjugate gradients results in optimal complexity.

I will also discuss the solution of diffusion-type problems
using exponential integrators based on Krylov subspace methods.

Numerical results will be presented for both regular and
scale-free graph topologies.

This is joint work with Mario Arioli (University of Wuppertal).


\end{document}
